% !TeX root = ../main.tex
\section{Finding settings: prediction and search}
\label{sec:search}

\subsection{Two ways to invert a chain}

\begin{description}[style=nextline,leftmargin=1.2cm]
\item[Prediction] Train a network on many renders of the chain with random
  settings to output the settings from the audio. Fast at use time (one pass),
  but it only knows the kind of audio it was trained on, and it struggles when
  several settings sound alike (it averages them).
\item[Analysis by synthesis (search)] Propose settings, render the dry signal
  through the real chain, measure the distance to the target, and improve the
  proposal; repeat. Slow (hundreds of renders), but it uses the exact forward
  model and needs no training.
\end{description}
The literature agrees on combining them: a prediction gives starting points, a
search refines them (Sections~\ref{sec:recognition} and~\ref{sec:matching}).
The app currently implements the search, starting from neutral settings and
from the factory presets.

\subsection{Gradient or no gradient}

If the chain is written in a differentiable framework (DDSP, PyTorch), the
gradient of the distance with respect to every setting is available and
gradient descent can be used. The app's chain is Rust code working on 16-bit
samples, with switches (pedal on or off) and discrete choices (amp voicing): it
has no gradients. So the search must be \term{gradient-free}: it may only
evaluate the distance at chosen points. That is also the situation of the
commercial plugins ST-ITO and Okita and Katayose search over
(Section~\ref{sec:recognition}). On non-convex problems gradient-free search can
even win: in synthesizer matching, CMA-ES reached a loss of 2.09 where Adam
(a gradient method) reached 4.60 \reported~\cite{bogdan2026instrumental}.

\subsection{CMA-ES, step by step}

\begin{keyidea}[: a cloud that moves, stretches and shrinks]
CMA-ES (Covariance Matrix Adaptation Evolution Strategy)~\cite{hansen2001cmaes}
keeps a Gaussian ``cloud'' over the search space: a \emph{mean} (the current
best guess), a \emph{step size} $\sigma$ (how wide the cloud is) and a
\emph{covariance} (its shape: which directions are stretched). Each generation
it samples candidates from the cloud, evaluates them, moves the mean toward the
best ones, stretches the cloud along directions that worked, and widens or
narrows it depending on whether the steps have been consistently in one
direction (far from the optimum: widen) or random (near it: narrow).
\end{keyidea}

One generation, with $n$ parameters, population $\lambda$ and the best
$\mu = \lambda/2$ kept:
\begin{enumerate}
\item \textbf{Sample} $\lambda$ candidates $x_k = m + \sigma\, C^{1/2} z_k$,
  with $z_k$ standard normal.
\item \textbf{Evaluate} each candidate (here: render and compare). This is the
  expensive part and runs in parallel.
\item \textbf{Rank} them and form the new mean as a weighted average of the
  best $\mu$, with larger weights for better ranks:
  $m' = \sum_{i=1}^{\mu} w_i\, x_{(i)}$.
\item \textbf{Update the shape} $C$ from the steps of the best candidates
  (rank-$\mu$ update) and from the accumulated path of the mean (rank-one
  update), so directions of progress grow.
\item \textbf{Update the step size}: compare the length of the accumulated path
  of the mean with the length expected of a random walk; longer means the
  search is moving consistently, so $\sigma$ grows; shorter means it is
  circling the optimum, so $\sigma$ shrinks.
\end{enumerate}
Only the \emph{ranks} of the candidates matter, not the distance values, which
makes CMA-ES robust to the scale and noise of the objective. The
\term{separable} variant~\cite{ros2008sepcmaes} keeps only the diagonal of $C$
(one width per parameter, no correlations): cheaper and faster to adapt, and
adequate for about ten parameters. The app uses it (\file{tone/cmaes.rs}).

\begin{worked}[: budget]
The app runs three searches (one per amp voicing), each with a population of
16 for 30 generations: $3 \times (1 + 16\cdot30) = 1443$ renders, plus one
render per factory preset for the baseline. Each render processes a 12\,s
excerpt through the chain. In the committed benchmark one estimate took about
15\,s on 28 threads \measured.
\end{worked}

\subsection{Pitfalls of search}

\begin{itemize}
\item \textbf{Rugged landscapes}: LFO rates and delay times
  (Section~\ref{sec:dsp}). Okita and Katayose had to fix the chorus rate because
  its landscape was too rugged \reported~\cite{okita2026audio}. Remedy: measure
  those parameters from the signal instead.
\item \textbf{Discrete choices} (amp voicing, which capture, which IR) are not
  continuous: search each one separately, or treat them as candidates.
\item \textbf{The objective can be gamed}: TAPE showed a search driving an
  embedding distance down while the spectral error rose above that of the
  unprocessed input \reported~\cite{long2026tape}. Use an objective close to
  what listeners hear, and check the result with a different score.
\item \textbf{Ambiguity}: many settings reach the same distance. A
  \term{prior} over realistic presets (a probability distribution fitted on
  real presets) turns the search into a maximum-a-posteriori estimate and
  reduced parameter error by up to a third in Yu et al.\ \reported~\cite{yu2025improving}.
\end{itemize}
